\scanpage{29}
$\unsure{C^{\bullet}|C^{\cdot}}(A'/A)=\varprojlim_n C^{\bullet}(A'_n/A_n)$, donc en vertu de M.L. on
aura
\[
\begin{aligned}
H^1(A'/A)&=H^1(C^{\bullet}(A'/A))=\varprojlim H^1(C^{\bullet}(A'_n/A_n))\\
&=\varprojlim_n H^1(A'_n/A_n)=0
\end{aligned}
\]

On peut aussi déduire formellement (i) de
(ii) sans difficulté (en faisant $m=1$, et $m=2$)

(ii) Va résulter de (iii), par le $\overline{\mathrm m}$ argument que
dans a).

(iii) \deleted{\unsure{So}} On voit comme dans (i) qu'on peut
supposer $A$ \deleted{\uline{local}, \unsure{donc}} \uline{complet}, \deleted{et évidemment
aussi \uline{local}. Soit Enfin, on} On peut aussi supposer
$A=B$, et $A$ \uline{local}. Soit $E'_{\deleted{\unsure{A}}}=E_{A'}$ plat de
type fini sur $A'$, avec donnée de descente, et
soit $E=H^0_D(A'/A,E')$. Il faut prouver
\unsure{au|que} 1°) $E$ plat sur $A$ \qquad 2°) $E\otimes_A A'\to E'$ est un isomorphisme. Or on aura
\[
C^{\bullet}_D(A'/A,E')\simeq\varprojlim_n C^{\bullet}_D(A'_n/A_n,E'_n)
\]
d'où \deleted{\unsure{$E$}} par M.L.
\[
E=H^0_D(A'/A,E')=\varprojlim_n H^0(A'_n/A_n,E'_n)
\]
Posons
\[
E^{(n)}=H^0(A'_n/A_n,E'_n)
\]
alors grâce à \deleted{\unsure{a)}} l'hypothèse a), sous la forme (iii),
on voit que $E^{(n)}$ est un $B^{(n)}$-module plat de
T.f., et $E^{(n)}\otimes_{B^{(n)}} A'_n\to E'_n$ est un isom.
% archive sheet number: 28
